\section{Corporate Actions and Adjusted Prices}

Stock splits and dividends are retained as vendor events. Representative
split-event windows are audited for continuity in adjusted returns and for
changes in reported share volume. Signed adjustment factors are not inferred
from the direction of a price change alone. The volume adjustment policy is
documented explicitly; reported share volume is not relabelled as dollar
volume, and normalization for splits is not assumed without evidence.

\InputGenerated{generated/corporate_action_results.tex}

\IfFileExists{generated/table_action_counts.tex}{
  \begin{table}[htbp]
    \centering
    \caption{Recorded dividend and split events by year.}
    \label{tab:action_counts}
    \small
    \begin{tabular}{@{}rrr@{}}
\toprule
Year & Split events & Dividend events \\
\midrule
2016 & 100 & 4,765 \\
2017 & 99 & 4,905 \\
2018 & 91 & 5,067 \\
2019 & 104 & 5,221 \\
2020 & 98 & 4,854 \\
2021 & 77 & 5,150 \\
2022 & 116 & 5,523 \\
2023 & 198 & 5,618 \\
2024 & 225 & 5,623 \\
2025 & 205 & 5,656 \\
\bottomrule
\end{tabular}

  \end{table}
}{\PendingResults{generated/table_action_counts.tex}}

\FigureGenerated{../figures/corporate_actions_by_year.pdf}
  {Recorded dividend and stock-split events in the validated Yahoo histories.}
  {fig:action_counts}

\IfFileExists{generated/table_dividend_audit.tex}{
  \begin{table}[htbp]
    \centering
    \caption{Representative dividend-event continuity checks; cash amounts are USD per share.}
    \label{tab:dividend_audit}
    \small
    \begin{tabular}{@{}llrrrrl@{}}
\toprule
Security & Ex-date & Cash amount & Adj. return & Raw total return & Abs. error & Status \\
\midrule
ABM & 2016-01-05 & 0.1650 & 0.236\% & 0.235\% & 0.0014\% & pass \\
ACU & 2016-01-05 & 0.1000 & 0.817\% & 0.812\% & 0.0047\% & pass \\
DHR & 2016-07-05 & 24.5600 & 22.228\% & 5.653\% & 16.5752\% & review \\
ABM & 2017-01-03 & 0.1700 & 0.320\% & 0.318\% & 0.0013\% & pass \\
BRX & 2017-01-03 & 0.2600 & 1.407\% & 1.392\% & 0.0150\% & pass \\
APLE & 2018-01-02 & 0.1000 & 0.461\% & 0.459\% & 0.0023\% & pass \\
CMCSA & 2018-01-02 & 0.1481 & 2.953\% & 2.941\% & 0.0116\% & pass \\
KDP & 2018-07-10 & 103.7500 & 11.451\% & 1.844\% & 9.6077\% & review \\
PPSI & 2019-10-08 & 1.3700 & 29.003\% & 20.513\% & 8.4902\% & review \\
TRTX & 2020-03-26 & 0.4300 & 92.953\% & 81.232\% & 11.7213\% & review \\
VHC & 2023-04-18 & 20.0000 & 21.154\% & 7.237\% & 13.9170\% & review \\
RBCN & 2023-10-24 & 1.1000 & 13.636\% & 5.114\% & 8.5227\% & review \\
\bottomrule
\end{tabular}

  \end{table}
}{\PendingResults{generated/table_dividend_audit.tex}}

\IfFileExists{generated/table_split_audit.tex}{
  \begin{table}[htbp]
    \centering
    \caption{Representative split-event audit windows.}
    \label{tab:split_audit}
    \small
    \begin{tabular}{@{}llrrrrl@{}}
\toprule
Security & Event date & Split & Adj. return & Raw return & Volume ratio & Status \\
\midrule
SHEN & 2016-01-05 & 2 & -0.10\% & -0.10\% & 0.48 & pass \\
DHR & 2016-07-05 & 1.319 & 22.23\% & -21.42\% & 3.75 & review \\
RILY & 2016-11-25 & 1.01 & 2.52\% & 1.00\% & 0.27 & pass \\
CFBK & 2018-08-21 & 0.181 & -1.70\% & -1.70\% & 0.60 & pass \\
CMCT & 2019-09-03 & 0.3333333333333333 & -8.78\% & -72.63\% & 0.85 & review \\
CTO & 2020-11-18 & 1.228 & 0.88\% & -1.92\% & 0.72 & pass \\
MRK & 2021-06-03 & 1.048 & 2.09\% & 2.09\% & 1.35 & pass \\
FWONA & 2023-07-20 & 1.018 & -1.07\% & -1.07\% & 2.76 & pass \\
CZFS & 2024-06-14 & 1.01 & 0.24\% & -0.90\% & 0.88 & pass \\
BBGI & 2024-09-24 & 0.05 & 5.52\% & 5.52\% & 1.11 & pass \\
CRESY & 2025-11-28 & 1.025 & 4.96\% & -0.50\% & 0.99 & review \\
RPT & 2025-12-31 & 0.16666666666666666 & 1.22\% & 1.22\% & 1.18 & pass \\
\bottomrule
\end{tabular}

  \end{table}
}{\PendingResults{generated/table_split_audit.tex}}

\FigureGenerated{../figures/split_adjustment_audit.pdf}
  {Representative before-and-after adjusted-price continuity around splits.}
  {fig:split_audit}

\subsection{Corporate-action reconciliation}

The prior validation's three split-continuity flags were independently
reproduced and all coincide with a distribution on the same date. The current
vendor-aware validator moves those cases to the combined-action issue category,
leaving \SplitFlags{} simple split-continuity flags. The dividend flags were
reproduced and classified into special distributions, large adjustment events,
simultaneous split/distribution events, and ordinary dividend mismatches.

\IfFileExists{generated/table_corporate_action_reconciliation.tex}{
  \begin{table}[htbp]
    \centering
    \caption{Representative corporate-action reconciliation.}
    \label{tab:corporate_action_reconciliation}
    \scriptsize
    \resizebox{\linewidth}{!}{{\small
\begin{tabular}{@{}llp{0.30\linewidth}p{0.24\linewidth}ll@{}}
\toprule
Security & Date & Action & Classification & confidence & Status \\
\midrule
DHR & 2016-07-05 & split 1.319; distribution 24.56 & simultaneous split and distribution & high & resolved event type \\
CMCT & 2019-09-03 & split 0.3333; distribution 102,941.18 & simultaneous split and distribution & high & resolved methodology \\
CRESY & 2025-11-28 & split 1.025; distribution 0.62 & simultaneous split and distribution & high & resolved methodology \\
KDP & 2018-07-10 & distribution 103.75 & probable special distribution or spinoff & high & resolved event type \\
RBCN & 2022-08-30 & distribution 11.00 & probable special distribution or spinoff & medium & partially resolved \\
VHC & 2023-04-18 & distribution 20.00 & probable special distribution or spinoff & medium & partially resolved \\
RBCN & 2023-10-24 & distribution 1.10 & probable special distribution or spinoff & medium & partially resolved \\
MBI & 2023-12-26 & distribution 8.00 & probable special distribution or spinoff & medium & partially resolved \\
CLRO & 2023-06-01 & distribution 15.00 & probable special distribution or spinoff & medium & partially resolved \\
GEN & 2020-02-03 & distribution 12.00 & probable special distribution or spinoff & medium & partially resolved \\
VC & 2016-01-25 & distribution 43.40 & probable special distribution or spinoff & medium & partially resolved \\
\bottomrule
\end{tabular}
}
}
  \end{table}
}{\PendingResults{generated/table_corporate_action_reconciliation.tex}}

\FigureGenerated{../figures/investigation/anomaly_corporate_action_reconciliation.pdf}
  {Original corporate-action flags and reconciled evidence classifications.}
  {fig:anomaly_corporate_actions}

Danaher's 2016 separation and distribution is independently documented in its
SEC Form 8-K \citep{secDHR2016}. Keurig Dr Pepper's 2018 special cash dividend
is independently documented in its SEC Form 8-K \citep{secKDP2018}. These
resolve the event type but not every detail of Yahoo's adjustment arithmetic.

\textbf{Evidence and uncertainty.} The tests recomputed raw and adjusted gross
returns, dividend fractions, split metadata, and vendor repair flags on each
flagged event. The evidence supports the conclusion that the original split flags are
combined actions and that most dividend flags are distributions rather than
ordinary dividends. Most non-special events lack an independently verified
event-type source, and Yahoo does not expose structured distribution
classifications. The recommended action is an action-specific continuity test,
not an assumption that every observed distribution belongs in the same formula.
